% Figure Legends (250 words per legend)
% - Should always have a short, explanatory title (as well as legend).
% - Legends must explain the figure fully without reference to the text.
% - If applicable, the original source of previously published material must be acknowledged in the Figure legend.
% - Figures must be cited in the main text (box figures must be cited in the box). 

\begin{figure}[p]
  \centering
  \includegraphics[width=14cm]{figures/architecture}
  \caption{\textbf{Typical neural network components of a deep agent.}
    (a) A \emph{visual processing} module (typically a convolutional
    network) converting pictures into internal distributed
    representations.  (b) A \textit{generation} component consisting
    of a recurrent neural network that produces a symbol sequence
    (in this case, $AXZ$).  (c) An
    \textit{understanding} module, that takes as input a sequence of
    units (in this case, the symbols produced by the generation
    component) and produces an internal distributed representation.  A typical
    \textit{sender} agent will first transform images into distributed
    representations with (a) and then use (b) to produce a message.
    A \textit{receiver} agent will also use (a) to transform images to
    representations, and then (c) to process the message from the
    sender in order to make a decision about the output action. In
    both cases, further layers are interspersed with the
    various components to further aid the agents' ``reasoning''
    process (e.g., the receiver might use them to combine visual and verbal information).
    \label{fig:agents}}
\end{figure}
